% Figure: DC load line of a common-emitter stage. The transistor output
% characteristics (I_C vs V_CE for a family of base currents) are the rising
% curves; the load line V_CE = V_CC - I_C (R_C + R_E) is the straight line;
% the bias point sits where the chosen base-current curve crosses the load
% line. pgfplots. V_CC = 12 V, (R_C + R_E) = 2670 ohm, so the load line
% runs from (12, 0) to (0, 12/2670 = 4.49 mA).
\begin{figure}[htbp]
\centering
\begin{tikzpicture}
\begin{axis}[
  width=12cm, height=7.5cm,
  xlabel={collector-emitter voltage $V_{CE}$ (V)},
  ylabel={collector current $I_C$ (mA)},
  xmin=0, xmax=12.5, ymin=0, ymax=5,
  axis lines=left,
  tick label style={font=\scriptsize},
  label style={font=\small},
  legend pos=north east,
  legend style={font=\scriptsize},
  clip=false,
]
% Output characteristics: near-flat active-region curves, each labelled by I_B.
% Model I_C = beta I_B (1 + V_CE/V_A) with Early voltage V_A = 80 V, beta = 150.
% I_B values chosen so the curves bracket the bias point at I_C ~ 2.67 mA.
\addplot[enggray, thick, domain=0.2:12, samples=60] {150*0.010*(1 + x/80)};
\node[enggray, font=\scriptsize, anchor=west] at (axis cs:12,{150*0.010*(1+12/80)}) {$I_B{=}10\,\mu\mathrm{A}$};
\addplot[enggray, thick, domain=0.2:12, samples=60] {150*0.0155*(1 + x/80)};
\node[enggray, font=\scriptsize, anchor=west] at (axis cs:12,{150*0.0155*(1+12/80)}) {$I_B{=}15.5\,\mu\mathrm{A}$};
\addplot[enggray, thick, domain=0.2:12, samples=60] {150*0.021*(1 + x/80)};
\node[enggray, font=\scriptsize, anchor=west] at (axis cs:12,{150*0.021*(1+12/80)}) {$I_B{=}21\,\mu\mathrm{A}$};
% DC load line from (V_CC, 0) to (0, V_CC/(R_C+R_E))
\addplot[engblue, very thick, domain=0:12, samples=2] {(12 - x)/2.670};
\addlegendentry{load line $V_{CE}=V_{CC}-I_C(R_C{+}R_E)$}
% bias point (quiescent): V_CE = 4.9 V, I_C = 2.67 mA
\filldraw[engred] (axis cs:4.9,2.67) circle [radius=2.2pt];
\node[engred, font=\scriptsize, anchor=south west] at (axis cs:5.05,2.72) {$Q$ point $(4.9\,\mathrm{V},\,2.67\,\mathrm{mA})$};
% saturation and cutoff markers
\node[enggray, font=\scriptsize, anchor=south] at (axis cs:0.6,4.5) {saturation};
\node[enggray, font=\scriptsize, anchor=south east] at (axis cs:12,0.35) {cutoff};
\end{axis}
\end{tikzpicture}
\caption[Common-emitter DC load line and bias point]{The DC load line of a
common-emitter stage and its intersection with the transistor output
characteristics. The straight load line is fixed by the supply and the total
collector-plus-emitter resistance, running from the supply voltage on the
horizontal axis to the short-circuit collector current on the vertical axis.
The gently rising grey curves are the transistor's output characteristics, one
per base current, their small slope the Early-effect output conductance. The
quiescent point sits where the base current set by the divider crosses the load
line, here near the centre of the line, which gives the widest symmetric output
swing before the stage runs into saturation at the left or cutoff at the right.
Emitter degeneration flattens the load line's dependence on the base current,
which is how the bias point holds against the part-to-part scatter in current
gain.}
\label{fig:vol04ch09:loadline}
\end{figure}
